% figures/ch03-line-bundle.tex
% Schematic: a line bundle and a section on a compact Riemann surface
\begin{figure}[ht]
\centering
\begin{tikzpicture}[>=Stealth, font=\small]
  \fill[blue!8] (0,0) .. controls (1.2,1.4) and (3.8,1.5) .. (5.2,0.3)
                .. controls (4.5,-1.2) and (1.0,-1.2) .. (0,0);
  \draw[thick,blue!60!black] (0,0) .. controls (1.2,1.4) and (3.8,1.5) .. (5.2,0.3)
                .. controls (4.5,-1.2) and (1.0,-1.2) .. (0,0);
  \node at (2.5,-1.0) {$X(\Gamma)$};

  \foreach \x/\h in {0.8/1.3,1.7/1.6,2.6/1.4,3.5/1.8,4.4/1.2}
    \draw[gray!70] (\x,0.2) -- (\x,\h);

  \draw[thick,orange!80!black] (0.8,0.9) .. controls (1.6,1.15) and (2.5,1.0) .. (3.5,1.25)
                                .. controls (4.0,1.35) and (4.2,1.0) .. (4.4,0.9);
  \fill[orange!80!black] (1.7,1.15) circle (1.8pt);
  \fill[orange!80!black] (3.5,1.25) circle (1.8pt);

  \node[above] at (1.7,1.6) {纤维};
  \node[orange!80!black,above] at (3.2,1.45) {截面 $s_f$};
  \node[left] at (0.8,0.2) {$P$};
  \node[right] at (4.4,0.2) {$Q$};

  \draw[decorate,decoration={brace,amplitude=5pt}] (5.6,0.2) -- (5.6,1.8)
    node[midway,right=6pt] {$\mathcal L^{k/2}$ 的局部样子};
\end{tikzpicture}
\caption{把权 $k$ 模形式看成线丛 $\mathcal L^{k/2}$ 的截面。权 $2$ 时这条线丛就是规范丛 $\Omega^1$；更高权只是把这条线丛张量幂化。Riemann--Roch 的任务就是：根据线丛的次数，计算它有多少条线性独立的截面。}
\label{fig:line-bundle}
\end{figure}
